% ---
% title: Bank products
% tags: finance, inclusion-exclusion, complement
% ---

\begin{problem}
	\begin{context}
		Banks try to sell additional products to existing clients
		(\emph{cross-selling}), because an existing client is cheaper to reach than
		a new one. Knowing who already has which product tells the bank whom to
		target.
	\end{context}
	A bank has 1000 retail clients: 820 have a current account ($C$), 350 a
	credit card ($K$), 240 a mortgage ($M$); 300 have both $C$ and $K$, 200
	have $C$ and $M$, 90 have $K$ and $M$, and 70 have all three. (a) How many
	clients have none of the products? (b) How many have exactly one product?
	(c) A campaign targets clients with a current account but no credit card.
	How many are they?
\end{problem}

\begin{slides}
	\small
	\exstatement
	\pause
	\begin{columns}[T]
		\begin{column}{.68\textwidth}
			\sol
			(a) Inclusion--exclusion:
			$|C\cup K\cup M| = 820+350+240-300-200-90+70 = 890$; so
			$1000-890=110$ clients have none.\\[.5mm]
			(b) Fill the diagram from the inside out: center $70$; then
			$300-70=230$, $200-70=130$, $90-70=20$; then only $C$:
			$820-230-130-70=390$, only $K$: $30$, only $M$: $20$; total
			$440$.\\[.5mm]
			(c) $|CK^c| = |C|-|CK| = 820-300 = 520$ (in the diagram: $390+130$).
		\end{column}
		\begin{column}{.3\textwidth}
			\centering
			\begin{tikzpicture}[scale=.52]
				\draw[thick] (-2.2,-3.6) rectangle (4,2.1);
				\draw[navy,thick] (0,0) circle (1.5);
				\draw[teal,thick] (1.8,0) circle (1.5);
				\draw[gold,thick] (.9,-1.5) circle (1.5);
				\node at (-.7,.3) {\scriptsize 390};
				\node at (2.5,.3) {\scriptsize 30};
				\node at (.9,-2.3) {\scriptsize 20};
				\node at (.9,.6) {\scriptsize 230};
				\node at (.05,-.95) {\scriptsize 130};
				\node at (1.75,-.95) {\scriptsize 20};
				\node at (.9,-.45) {\scriptsize 70};
				\node at (3.3,-3.1) {\scriptsize 110};
				\node[navy] at (-1.4,1.4) {$C$};
				\node[teal] at (3.2,1.4) {$K$};
				\node[gold] at (-.9,-2.7) {$M$};
				\node at (.9,1.8) {\scriptsize $\Omega$};
			\end{tikzpicture}
		\end{column}
	\end{columns}
	\pause
	\smallskip
	\alert{Next week:} divide by 1000 and these become probabilities, with
	$P(E\cup F)=P(E)+P(F)-P(EF)$.
\end{slides}

\begin{solution}
	\idea Inclusion--exclusion for three sets, and ``from the inside out'' for the
	``only'' regions.

	\begin{minipage}{.56\linewidth}
		\step{Filling the Venn diagram.}
		\begin{itemize}[nosep]
			\item center (all three): $70$;
			\item only $C$ and $K$: $300-70=230$;\\ only $C$ and $M$:
			      $200-70=130$;\\ only $K$ and $M$: $90-70=20$;
			\item only $C$: $820-230-130-70=390$;\\ only $K$:
			      $350-230-20-70=30$;\\ only $M$: $240-130-20-70=20$;
			\item outside: $1000-890=110$ (part (a)).
		\end{itemize}
	\end{minipage}\hfill
	\begin{minipage}{.4\linewidth}
		\centering
		\begin{tikzpicture}[scale=.65]
				\draw[thick] (-2.2,-3.6) rectangle (4,2.1);
				\draw[navy,thick] (0,0) circle (1.5);
				\draw[teal,thick] (1.8,0) circle (1.5);
				\draw[gold,thick] (.9,-1.5) circle (1.5);
				\node at (-.7,.3) {\small 390};
				\node at (2.5,.3) {\small 30};
				\node at (.9,-2.3) {\small 20};
				\node at (.9,.6) {\small 230};
				\node at (.05,-.95) {\small 130};
				\node at (1.75,-.95) {\small 20};
				\node at (.9,-.45) {\small 70};
				\node at (3.3,-3.1) {\small 110};
				\node[navy] at (-1.4,1.4) {$C$};
				\node[teal] at (3.2,1.4) {$K$};
				\node[gold] at (-.9,-2.7) {$M$};
				\node at (.9,1.8) {\small $\Omega$};
			\end{tikzpicture}
	\end{minipage}

	\step{(a)}
	\[ |C\cup K\cup M|=820+350+240-300-200-90+70=890. \]
	None: $1000-890=110$. Check: the seven regions of the diagram sum to
	$390+30+20+230+130+20+70=890$.

	\step{(b)} ``Only $C$'' $=|C|-|CK|-|CM|+|CKM|$: removing $CK$ and $CM$
	removes the 70 clients with all three products twice, so we add them back
	once.
	\[ \begin{aligned}
			\text{Only }C & \colon\ 820-300-200+70=390; \\
			\text{only }K & \colon\ 350-300-90+70=30;   \\
			\text{only }M & \colon\ 240-200-90+70=20.
		\end{aligned} \]
	Exactly one product: $390+30+20=440$, the three outer regions of the
	diagram.

	\step{(c)} $C$ splits into $CK$ (has a card) and $CK^c$ (no card):
	\[ |CK^c|=|C|-|CK|=820-300=520. \]
	In the diagram these are the two regions of $C$ outside $K$: $390+130$.
	\begin{answer}
		(a) $110$ \quad (b) $440$ \quad (c) $520$
	\end{answer}
	\begin{remark}
		Dividing by 1000 turns these counts into probabilities for a randomly
		chosen client; e.g.\ $P(C\cup K\cup M)=0.89$. This is exactly how
		probabilities will be computed next week.
	\end{remark}
\end{solution}
